\documentclass{zhvt-classic}

% 使用本地源流明体（GenRyuMin2 TW）按文件路径加载，不依赖系统字体名
% 版框（内框／界栏／书口侧栏／鱼尾）改用与版心印章同一印泥棕色：
%   印章图片 yshs.png 实色 = RGB(135,68,52) = HTML 874434
\usepackage{xcolor}
\usepackage{tikz}
% sealbrown 双注册：xcolor 供 TikZ/pgf，L3 供 cls 版框
\definecolor{sealbrown}{HTML}{874434}
\ExplSyntaxOn
\color_set:nnn { sealbrown } { HTML } { 874434 }
\ExplSyntaxOff

\zhvtset{
  grid color       = sealbrown,
  main font        = GenRyuMin2TW-SB,
  main font path   = C:/Users/Administrator/texmf/fonts/opentype/public/GenRyuMin2TW/,
jiazhu font      = GenRyuMin2TW-M,
  jiazhu font path = C:/Users/Administrator/texmf/fonts/opentype/public/GenRyuMin2TW/,
  % 页码框容纳 4 字：正常页「九十七」+前置空 1 格=4 行恰好；
  % 100+ 页「一百零一」共 5 行超出，末字「一/二/三」被框裁掉（2026-10-04 实测），故放宽到 5
  shukou page number chars = 5,
  % 书口书名框高从默认 6 字减到 4 字，把省出的 2 字高拨给 chapter 框，
  % 使 6 字书名+卷名放进 9 字高的 chapter 框后不再压线（2026-10-05）
  shukou title chars = 4,
}

\title{史记}[袁盎晁错列传]
\author{汉·司马迁}
\maker{闲叟漫记·重製}


% ===== 索隐述赞样式（2026-11-08 四次定稿 · 愚大人亲手改定参数）=====
% ⚠ 竖排前提：自定义字体家族必须带 RawFeature={vertical:+vert}；
%   simhei.ttf 加 +vhal 会整体加载失败静默回退，只写 +vert；simfang 同理。
\newCJKfontfamily{\shuzanfont}[Path=C:/Windows/Fonts/,RawFeature={vertical:+vert}]{simhei.ttf}
\newCJKfontfamily{\shuzanbody}[Path=C:/Windows/Fonts/,RawFeature={vertical:+vert}]{simfang.ttf}
% 题头：152x30pt 圆角 6pt 棕底白字；baseline (0,2) 使整框落入界行格子不压版框线；
% 白字节点 x 不补偿、y +8.2pt（像素级墨迹实测 dx/dy≈0，L/R 留白约 20pt）。
\newcommand{\shuzanhead}[1]{%
  \begin{tikzpicture}[baseline={(0,2)}]
    \fill[sealbrown, rounded corners=6pt]
      (-76pt,-15pt) rectangle (76pt,15pt);
    \node[text=white, inner sep=0pt, anchor=center] at (0pt,8.2pt)
      {\shuzanfont\fontsize{23}{23}\selectfont #1};
  \end{tikzpicture}%
}
% ★ 修正书口卷名位置：chapter 框起始位置改回 cls 默认值，避免全局影响。
%   只在每个卷首页通过 \juan 命令里的空命令微调该页位置（安全，不会跨页污染）。
\ExplSyntaxOn
\AtBeginDocument{
  \dim_set:Nn \l__zhvt_shukou_chapter_title_start_dim
    { \l__zhvt_shukou_title_height_dim + \l__zhvt_font_size_dim * 2 }
  \dim_set:Nn \l__zhvt_shukou_page_number_start_dim
    { \l__zhvt_shukou_chapter_title_start_dim + \l__zhvt_shukou_chapter_title_height_dim }
}
\ExplSyntaxOff

\begin{document}

% ★★ 版心（书口）排布覆写：把书名放到鱼尾之下（2026-10-04）
%
%   cls 默认版心自上而下：书名（title 静态框，在鱼尾之上）→ 鱼尾 → 篇名（chapter 动态框）→ 页码
%   本模板改为：          鱼尾（位置不动）→ 书名 → 篇名 → 页码分割线 → 页码
%
%   手法：把 cls 放在 title 静态框里的书名，改放进 chapter 动态框（鱼尾之下）；
%         title 静态框留空。书名取 \l__zhvt_title_str，即 \title{...} 的第一参数，
%         所以换书时只改 \title，不必在此重复写书名。
%
%   ⚠ 三条必须遵守的写法（2026-10-04 实测）：
%     ① 汉字必须用 \hbox{...} 包住。flowfram 会把框内容装进一个宽 14.4mm（= 一个 line_sep
%        = 书口栏宽）的纵向 minipage 并按该宽度折行，而整页靠 PDF /Rotate 90 显示、
%        框内文字本身不变向 —— 不包 \hbox 时「围炉夜话」4 字会被逐字折成 4 行，
%        在版面上就成了横排 4 字并横向溢出版框。
%     ② 竖排下**一个 \hbox 就是一整列**。书名与篇名要同处一列，必须放进同一个 \hbox；
%        分成两个 \hbox 时，后一个会被当成新的一列、甩到隔壁栏去
%        （实测「序」跑到书口栏外侧 14.4mm 的正文列上）。中间的 \quad 在框内是横向空白，
%        显示出来正好是书名与篇名之间的纵向间隔。
%     ③ ★ 内容盒要写成 **与框等宽的定宽盒** `\hbox to \l__zhvt_shukou_chapter_title_height_dim{…\hfil}`。
%        实测：chapter 框是把内容盒沿纵向**居中**摆放的，内容宽度一变，书名就跟着漂：
%        自然宽度的 \hbox（书名 38.4mm）落在 x=157.2mm，加篇名后（57.6mm）落在 x=152.4mm，
%        目录页（无篇名）与正文页的书名就差半格。做成等宽盒后，无论有没有篇名，
%        内容固定从框顶起排，实测两页均落在 x=142.8mm，不再漂。
\ExplSyntaxOn
% 把书口 chapter 框（书名+篇名）起始位置下移 1 字高，避免 4 字卷名顶到鱼尾
% ⚠ cls 的 \__zhvt_calculate_dimensions: 在 \AtBeginDocument 里执行，会重算该 dim；
%   必须把覆写挂到 begin document 之后，否则会被二次覆盖（2026-10-05 实测）。
% (旧方案已移到导言区，直接修改 chapter 框起始位置)
% 卷标题命令 #1=书口篇名(鱼尾) #2=正文卷名（含卷序，卷名内空格用 ~）
%
% ★★ 目录页两项修订（2026-10-05，愚大人指出）★★
%   ① 页码起点统一对齐：
%      原状 —— 目录项写成 \par\hspace{4em}卷名\hspace{4em}页码，页码起点 = 卷名实宽 + 4em。
%      卷名实宽随字数变：「卷一~梁惠王上」6 字 +~、「卷七~离娄上」5 字 +~，
%      差 1 个字位 → 卷七~卷十 4 条的页码整体比其余 10 条高 27.21pt（一个字），看着上下不齐。
%      改法 —— 卷名补全角空格（宽 = 1 字 = 27.215pt）补足到 **7 字符位**再写进目录，
%      14 条页码起点遂全部落在同一条竖线上（起点对齐；末字随页码字数自然参差）。
%      ⚠ 只对目录里的副本 \l_tmpc_tl 补，正文卷名仍用 \l_tmpa_tl 原样 —— 正文像素不受影响。
%   ② 卷名整体上移半字位：前导 \hspace 由 4em 改 3.5em。
%      实测原状目录块在版框内偏下（天头留白 139.6pt / 地脚留白 110.7pt），
%      上移 0.5em = 13.6pt 后上下基本均分，视觉居中。
\NewDocumentCommand { \juan } { mm } {
	\clearpage
	\tl_set:Nn \l__zhvt_chapter_title_tl { #1 }
	\tl_set:Nn \l_tmpa_tl { #2 }
	\tl_set_eq:NN \l_tmpc_tl \l_tmpa_tl
	\int_step_inline:nn { 7 - \str_count:n { #2 } }
	  { \tl_put_right:Nn \l_tmpc_tl { 　 } }
	\tl_set:Nx \l_tmpb_tl { \zhnum { page } }
	\tl_put_left:Nn \l_tmpb_tl { \hspace { 4em } }
	\tl_put_left:Nn \l_tmpb_tl { \l_tmpc_tl }
	\tl_put_left:Nn \l_tmpb_tl { \par \hspace { 3.5em } }
	\addtocontents { toc } { \l_tmpb_tl }
	\skip_horizontal:n { \l__zhvt_font_size_dim * 4 }
	\l_tmpa_tl
	\par
	\thispagestyle { empty }
}
\cs_set_protected:Npn \__zhvt_set_shukou_content: {
	% ★ 鱼尾上方留白（title 静态框）嵌入闲叟钤印：
	%   整页 /Rotate 90，故此处把图用 \rotatebox{90} 转正为竖章；
	%   尺寸卡在版心条幅内（书内约 7mm 宽 × 33.8mm 高，留余量不压鱼尾）。
	\setstaticcontents * { title } {
		\vbox_to_ht:nn { \l__zhvt_baseline_skip_dim } {
			\vskip 14.2pt
			\hbox_to_wd:nn { \l__zhvt_shukou_title_height_dim } {
				\hfil
				\rotatebox [ origin = c ] { 90 } { \includegraphics [ width = 7mm , height = 33.8mm ] { yshs.png } }
				\hfil
			}
		}
	}
% ★ 版心（书口）统一改法（2026-11-16 愚大人定）：18pt＋框内居中＋书口文字同版框色
%   ① 原 27.2pt 时 7 字长篇名（如「廉颇蔺相如列传」）末字溢出 chapter 框、压住页码；
%      改 \fontsize{18pt}{18pt}，前后 \hfil 令字串在框内居中（不随篇名字数漂移）。
%   ② \vbox{\vskip 8.2pt} 定量下移，抵消缩字号后栏宽方向 8.2pt 上偏（像素实测 y 中心 543.3 = 原版）。
%      ⚠ \raisebox / \vspace* 在 flowfram 框内无效；换字号须同步换补偿：16pt→10.1、19pt→7.4。
%   ③ chapter 与 page~number 均置 \color{sealbrown}（HTML 874434），与版框／界行同色；正文仍黑。
%   ⚠ 页码字位、钤印、鱼尾、界行、正文一律未动；五部合订本同此口径。
	\setdynamiccontents * { chapter }
	{ \color{sealbrown} \vbox { \vskip 8.2pt \hbox to \l__zhvt_shukou_chapter_title_height_dim
		{ \hfil \fontsize{18pt}{18pt}\selectfont \l__zhvt_title_str \quad \l__zhvt_chapter_title_tl \hfil } } }
	\setdynamiccontents * { page~number } { \color{sealbrown} \quad \zhnum { page } }
}
\__zhvt_set_shukou_content:
\ExplSyntaxOff

\maketitle

% 单篇，无多卷目录，不渲染目录页

\juan{袁盎晁错列传}{卷一百零一~袁盎晁错列传第四十一}
\hspace*{2em}袁盎者。楚人也。字丝。父故为群盗。徙处安陵。高后时。盎尝为吕禄舍人。及孝文帝即位。盎兄哙任盎为中郎。\\
\hspace*{2em}绛侯为丞相。朝罢趋出。意得甚。上礼之恭。常自送之。袁盎进曰。﹁陛下以丞相何如人。﹂上曰。﹁社稷臣。﹂盎曰。﹁绛侯所谓功臣。非社稷臣。社稷臣主在与在。主亡与亡。方吕后时。诸吕用事。擅相王。刘氏不绝如带。是时绛侯为太尉。主兵柄。弗能正。吕后崩。大臣相与共畔诸吕。太尉主兵。適会其成功。所谓功臣。非社稷臣。丞相如有骄主色。陛下谦让。臣主失礼。窃为陛下不取也。﹂后朝。上益庄。丞相益畏。已而绛侯望袁盎曰。﹁吾与而兄善。今兒廷毁我。﹂盎遂不谢。\\
\hspace*{2em}及绛侯免相之国。国人上书告以为反。徵系清室。宗室诸公莫敢为言。唯袁盎明绛侯无罪。绛侯得释。盎颇有力。绛侯乃大与盎结交。\\
\hspace*{2em}淮南厉王朝。杀辟阳侯。居处骄甚。袁盎谏曰。﹁诸侯大骄必生患。可適削地。﹂上弗用。淮南王益横。及棘蒲侯柴武太子谋反事觉。治。连淮南王。淮南王徵。上因迁之蜀。轞车传送。袁盎时为中郎将。乃谏曰。﹁陛下素骄淮南王。弗稍禁。以至此。今又暴摧折之。淮南王为人刚。如有遇雾露行道死。陛下竟为以天下之大弗能容。有杀弟之名。柰何。﹂上弗听。遂行之。\\
\hspace*{2em}淮南王至雍。病死。闻。上辍食。哭甚哀。盎入。顿首请罪。上曰。﹁以不用公言至此。﹂盎曰。﹁上自宽。此往事。岂可悔哉。且陛下有高世之行者三。此不足以毁名。﹂上曰。﹁吾高世行三者何事。﹂盎曰。﹁陛下居代时。太后尝病。三年。陛下不交睫。不解衣。汤药非陛下口所尝弗进。夫曾参以布衣犹难之。今陛下亲以王者脩之。过曾参孝远矣。夫诸吕用事。大臣专制。然陛下从代乘六传驰不测之渊。虽贲育之勇不及陛下。陛下至代邸。西向让天子位者再。南面让天子位者三。夫许由一让。而陛下五以天下让。过许由四矣。且陛下迁淮南王。欲以苦其志。使改过。有司卫不谨。故病死。﹂於是上乃解。曰。﹁将柰何。﹂盎曰。﹁淮南王有三子。唯在陛下耳。﹂於是文帝立其三子皆为王。盎由此名重朝廷。\\
\hspace*{2em}袁盎常引大体慷慨。宦者赵同以数幸。常害袁盎。袁盎患之。盎兄子种为常侍骑。持节夹乘。说盎曰。﹁君与斗。廷辱之。使其毁不用。﹂孝文帝出。赵同参乘。袁盎伏车前曰。﹁臣闻天子所与共六尺舆者。皆天下豪英。今汉虽乏人。陛下独奈何与刀锯馀人载。﹂於是上笑。下赵同。赵同泣下车。\\
\hspace*{2em}文帝从霸陵上。欲西驰下峻阪。袁盎骑。并车揽辔。上曰。﹁将军怯邪。﹂盎曰。﹁臣闻千金之子坐不垂堂。百金之子不骑衡。圣主不乘危而徼幸。今陛下骋六騑。驰下峻山。如有马惊车败。陛下纵自轻。柰高庙。太后何。﹂上乃止。\\
\hspace*{2em}上幸上林。皇后。慎夫人从。其在禁中。常同席坐。及坐。郎署长布席。袁盎引卻慎夫人坐。慎夫人怒。不肯坐。上亦怒。起。入禁中。盎因前说曰。﹁臣闻尊卑有序则上下和。今陛下既已立后。慎夫人乃妾。妾主岂可与同坐哉。適所以失尊卑矣。且陛下幸之。即厚赐之。陛下所以为慎夫人。適所以祸之。陛下独不见﹃人彘﹄乎。﹂於是上乃说。召语慎夫人。慎夫人赐盎金五十斤。\\
\hspace*{2em}然袁盎亦以数直谏。不得久居中。调为陇西都尉。仁爱士卒。士卒皆争为死。迁为齐相。徙为吴相。辞行。种谓盎曰。﹁吴王骄日久。国多奸。今苟欲劾治。彼不上书告君。即利剑刺君矣。南方卑湿。君能日饮。毋何。时说王曰毋反而已。如此幸得脱。﹂盎用种之计。吴王厚遇盎。\\
\hspace*{2em}盎告归。道逢丞相申屠嘉。下车拜谒。丞相从车上谢袁盎。袁盎还。愧其吏。乃之丞相舍上谒。求见丞相。丞相良久而见之。盎因跪曰。﹁原请间。﹂丞相曰。﹁使君所言公事。之曹与长史掾议。吾且奏之。即私邪。吾不受私语。﹂袁盎即跪说曰。﹁君为丞相。自度孰与陈平。绛侯。﹂丞相曰。﹁吾不如。﹂袁盎曰。﹁善。君即自谓不如。夫陈平。绛侯辅翼高帝。定天下。为将相。而诛诸吕。存刘氏。君乃为材官蹶张。迁为队率。积功至淮阳守。非有奇计攻城野战之功。且陛下从代来。每朝。郎官上书疏。未尝不止辇受其言。言不可用置之。言可受采之。未尝不称善。何也。则欲以致天下贤士大夫。上日闻所不闻。明所不知。日益圣智。君今自闭钳天下之口而日益愚。夫以圣主责愚相。君受祸不久矣。﹂丞相乃再拜曰。﹁嘉鄙野人。乃不知。将军幸教。﹂引入与坐。为上客。\\
\hspace*{2em}盎素不好晁错。晁错所居坐。盎去。盎坐。错亦去。两人未尝同堂语。及孝文帝崩。孝景帝即位。晁错为御史大夫。使吏案袁盎受吴王财物。抵罪。诏赦以为庶人。\\
\hspace*{2em}吴楚反。闻。晁错谓丞史曰。﹁夫袁盎多受吴王金钱。专为蔽匿。言不反。今果反。欲请治盎宜知计谋。﹂丞史曰。﹁事未发。治之有绝。今兵西乡。治之何益。且袁盎不宜有谋。﹂晁错犹与未决。人有告袁盎者。袁盎恐。夜见窦婴。为言吴所以反者。原至上前口对状。窦婴入言上。上乃召袁盎入见。晁错在前。及盎请辟人赐间。错去。固恨甚。袁盎具言吴所以反状。以错故。独急斩错以谢吴。吴兵乃可罢。其语具在吴事中。使袁盎为太常。窦婴为大将军。两人素相与善。逮吴反。诸陵长者长安中贤大夫争附两人。车随者日数百乘。\\
\hspace*{2em}及晁错已诛。袁盎以太常使吴。吴王欲使将。不肯。欲杀之。使一都尉以五百人围守盎军中。袁盎自其为吴相时。有从史尝盗爱盎侍兒。盎知之。弗泄。遇之如故。人有告从史。言﹁君知尔与侍者通﹂。乃亡归。袁盎驱自追之。遂以侍者赐之。复为从史。及袁盎使吴见守。从史適为守盎校尉司马。乃悉以其装赍置二石醇醪。会天寒。士卒饥渴。饮酒醉。西南陬卒皆卧。司马夜引袁盎起。曰。﹁君可以去矣。吴王期旦日斩君。﹂盎弗信。曰。﹁公何为者。﹂司马曰。﹁臣故为从史盗君侍兒者。﹂盎乃惊谢曰。﹁公幸有亲。吾不足以累公。﹂司马曰。﹁君弟去。臣亦且亡。辟吾亲。君何患。﹂乃以刀决张。道从醉卒隧出。司马与分背。袁盎解节毛怀之。杖。步行七八里。明。见梁骑。骑驰去。遂归报。\\
\hspace*{2em}吴楚已破。上更以元王子平陆侯礼为楚王。袁盎为楚相。尝上书有所言。不用。袁盎病免居家。与闾里浮沈。相随行。斗鸡走狗。雒阳剧孟尝过袁盎。盎善待之。安陵富人有谓盎曰。﹁吾闻剧孟博徒。将军何自通之。﹂盎曰。﹁剧孟虽博徒。然母死。客送葬车千馀乘。此亦有过人者。且缓急人所有。夫一旦有急叩门。不以亲为解。不以存亡为辞。天下所望者。独季心。剧孟耳。今公常从数骑。一旦有缓急。宁足恃乎。﹂骂富人。弗与通。诸公闻之。皆多袁盎。\\
\hspace*{2em}袁盎虽家居。景帝时时使人问筹策。梁王欲求为嗣。袁盎进说。其後语塞。梁王以此怨盎。曾使人刺盎。刺者至关中。问袁盎。诸君誉之皆不容口。乃见袁盎曰。﹁臣受梁王金来刺君。君长者。不忍刺君。然後刺君者十馀曹。备之。﹂袁盎心不乐。家又多怪。乃之棓生所问占。还。梁刺客後曹辈果遮刺杀盎安陵郭门外。\\
\hspace*{2em}晁错者。颍川人也。学申商刑名於轵张恢先所。与雒阳宋孟及刘礼同师。以文学为太常掌故。\\
\hspace*{2em}错为人穞直刻深。孝文帝时。天下无治尚书者。独闻济南伏生故秦博士。治尚书。年九十馀。老不可徵。乃诏太常使人往受之。太常遣错受尚书伏生所。还。因上便宜事。以书称说。诏以为太子舍人。门大夫。家令。以其辩得幸太子。太子家号曰﹁智囊﹂。数上书孝文时。言削诸侯事。及法令可更定者。书数十上。孝文不听。然奇其材。迁为中大夫。当是时。太子善错计策。袁盎诸大功臣多不好错。\\
\hspace*{2em}景帝即位。以错为内史。错常数请间言事。辄听。宠幸倾九卿。法令多所更定。丞相申屠嘉心弗便。力未有以伤。内史府居太上庙壖中。门东出。不便。错乃穿两门南出。凿庙壖垣。丞相嘉闻。大怒。欲因此过为奏请诛错。错闻之。即夜请间。具为上言之。丞相奏事。因言错擅凿庙垣为门。请下廷尉诛。上曰。﹁此非庙垣。乃壖中垣。不致於法。﹂丞相谢。罢朝。怒谓长史曰。﹁吾当先斩以闻。乃先请。为兒所卖。固误。﹂丞相遂发病死。错以此愈贵。\\
\hspace*{2em}迁为御史大夫。请诸侯之罪过。削其地。收其枝郡。奏上。上令公卿列侯宗室集议。莫敢难。独窦婴争之。由此与错有卻。错所更令三十章。诸侯皆諠譁疾晁错。错父闻之。从颍川来。谓错曰。﹁上初即位。公为政用事。侵削诸侯。别疏人骨肉。人口议多怨公者。何也。﹂晁错曰。﹁固也。不如此。天子不尊。宗庙不安。﹂错父曰。﹁刘氏安矣。而晁氏危矣。吾去公归矣。﹂遂饮药死。曰。﹁吾不忍见祸及吾身。﹂死十馀日。吴楚七国果反。以诛错为名。及窦婴。袁盎进说。上令晁错衣朝衣斩东市。\\
\hspace*{2em}晁错已死。谒者仆射邓公为校尉。击吴楚军为将。还。上书言军事。谒见上。上问曰。﹁道军所来。闻晁错死。吴楚罢不。﹂邓公曰。﹁吴王为反数十年矣。发怒削地。以诛错为名。其意非在错也。且臣恐天下之士噤口。不敢复言也。﹂上曰。﹁何哉。﹂邓公曰。﹁夫晁错患诸侯彊大不可制。故请削地以尊京师。万世之利也。计画始行。卒受大戮。内杜忠臣之口。外为诸侯报仇。臣窃为陛下不取也。﹂於是景帝默然良久。曰。﹁公言善。吾亦恨之。﹂乃拜邓公为城阳中尉。\\
\hspace*{2em}邓公。成固人也。多奇计。建元中。上招贤良。公卿言邓公。时邓公免。起家为九卿。一年。复谢病免归。其子章以脩黄老言显於诸公间。\\
\hspace*{2em}太史公曰。袁盎虽不好学。亦善傅会。仁心为质。引义慷慨。遭孝文初立。资適逢世。时以变易。及吴楚一说。说虽行哉。然复不遂。好声矜贤。竟以名败。晁错为家令时。数言事不用。後擅权。多所变更。诸侯发难。不急匡救。欲报私雠。反以亡躯。语曰﹁变古乱常。不死则亡﹂。岂错等谓邪。\\

% ===== 索隐述赞（司马贞《史记索隐》）· 统一样式 =====
% 题头「索隐述赞曰」五字并入反白框；正文每四句一列，定宽盒 \hfil 居列中；
% \fontsize{24}{41} 列距对齐全书界行（121+41k）；\clearpage 使述赞独占一页。
\clearpage
\shuzanhead{索隐述赞曰}\linebreak
{\shuzanbody\fontsize{24}{41}\selectfont
\hbox to \hsize{\hfil 袁丝公直。亦多附会。揽辔见重。卻席翳赖。\hfil}\linebreak
\hbox to \hsize{\hfil 晁错建策。屡陈利害。尊主卑臣。家危国泰。\hfil}\linebreak
\hbox to \hsize{\hfil 悲彼二子。名立身败。\hfil}%
\par}
\end{document}
